\documentclass[11pt,letterpaper]{article}
\usepackage{fontspec}
\usepackage{tocloft}
\usepackage{multicol}
\usepackage{nicefrac}
\usepackage[left=1.4in,right=1.5in,top=1in,bottom=1in]{geometry}
\setmainfont[Scale=1.4,AutoFakeBold=1.5,AutoFakeSlant=0.3]{Doves Type}

\widowpenalty=10000
\clubpenalty=10000
\interlinepenalty=500

\title{Restaurant-Style Salsa Roja}
\author{}
\date{}

\begin{document}

\maketitle
\thispagestyle{empty}
\vspace{-60pt}

% Intentionally suppressed
% \begin{quote}
% \textit{Canned whole tomatoes are pulsed with white onion, garlic, jalapeños, serrano, cilantro, cumin, lime, and salt into a bright, restaurant-style salsa. Ready in minutes; best after a short rest. Makes about 3~cups.}
% \end{quote}

\section*{Ingredients}
\setlength{\columnsep}{20pt}
\begin{multicols}{2}
\noindent
    Whole peeled tomatoes \dotfill 28 oz. can \\
    White onion \dotfill \nicefrac{1}{2} medium \\
    Garlic cloves \dotfill 5 \\
    Jalapeños \dotfill 4 \\
    \columnbreak
    Serrano chile \dotfill 1 \\
    Fresh cilantro \dotfill \nicefrac{1}{3} cup \\
    Lime juice \dotfill from \nicefrac{1}{2} lime \\
    Ground cumin \dotfill 1 tsp. \\
    Kosher salt \dotfill 1 Tbsp.
\end{multicols}

{\small\textit{*San~Marzano or other whole peeled tomatoes, undrained. Strip seeds and ribs from the chiles for milder heat.*}}

\section*{Directions}

\noindent
Roughly chop the \textbf{onion}, \textbf{garlic}, \textbf{jalapeños}, \textbf{serrano}, and \textbf{cilantro} --- Juice \nicefrac{1}{2}~\textbf{lime} --- Measure 1~tsp. \textbf{cumin} and 1~Tbsp. \textbf{kosher~salt}

\begin{enumerate}
    \item Add the \textbf{tomatoes} with their juice, chopped \textbf{onion}, \textbf{garlic}, \textbf{jalapeños}, \textbf{serrano}, \textbf{cilantro}, \textbf{lime~juice}, 1~tsp. \textbf{cumin}, and 1~Tbsp. \textbf{kosher~salt} to a food processor.

    \item Pulse in short bursts until the salsa reaches the texture you want, from chunky to nearly smooth. Stop and scrape the sides once or twice. The mixture should look even, with bits of chile and cilantro still visible. Do not run the processor continuously---that turns the salsa thin and foamy.

    \item Taste. The salsa should taste tomato-forward, spicy, and well salted, with a warm cumin note and a clean lime finish. Add more \textbf{salt} or \textbf{lime~juice} as needed. For thicker dip, pour off a splash of liquid from the can before blending next time, or drain a spoonful from the finished salsa.

    \item Rest \textit{15--30~minutes} so the flavors marry. Taste again and adjust \textbf{salt} if needed. Serve at cool room temperature. Best within a few hours; refrigerates up to \textit{5~days}. Stir before serving---it settles as it sits.
\end{enumerate}

\end{document}
